\documentclass[11pt]{article}
\usepackage[T1]{fontenc}
\usepackage[utf8]{inputenc}
\usepackage{amsmath,amssymb,amsthm}
\usepackage[margin=1in]{geometry}
\usepackage{hyperref}
\newtheorem{theorem}{Theorem}
\newtheorem{lemma}[theorem]{Lemma}
\newtheorem{corollary}[theorem]{Corollary}
\newcommand{\Hom}{\operatorname{Hom}}
\newcommand{\im}{\operatorname{im}}
\newcommand{\Fitt}{\operatorname{Fitt}}
\title{Faithfully flat descent of diagonal-compatible Cartier maps}
\author{Unreviewed mathematical working paper}
\date{8 October 2026}
\begin{document}
\maketitle
\begin{abstract}
For a finite type algebra over a perfect field, diagonal-compatible
Cartier maps and their restrictions commute with a perfect field
extension. In particular, liftability to the diagonal descends faithfully
flatly. The proof gives explicit Frobenius-module identifications and
uses finite presentations to commute kernels and images with extension.
It provides a replacement for a field-retraction argument in the finite
field step of a singular subadditivity source. Test-ideal base change,
the separate infinite-field inclusion, and the original envelope problem
are not asserted by this lemma.
\end{abstract}

\section{The precise modules and the theorem}
Let $k$ be perfect of characteristic $p>0$, let $R$ be a finite type
$k$-algebra, and fix $e\geq0$ and $q=p^e$. Neither reducedness nor
normality is needed for the module theorem. For a ring or ideal $M$
in characteristic $p$, $F_*^eM$ denotes its additive group with scalar
action $a\cdot F_*^ex=F_*^e(a^q x)$.
Put
\[
 S=R\otimes_k R,\quad \mu:S\longrightarrow R,\quad
 I=\ker\mu,\quad H_e(R)=\Hom_R(F_*^eR,R).
\]
Define the compatibility module
\[
 C_e(S,I)=\{\psi\in\Hom_S(F_*^eS,S):\psi(F_*^eI)\subset I\}.
\]
Its restriction to the quotient $S/I=R$ gives an $S$-linear map
\[
 \rho_e:C_e(S,I)\longrightarrow H_e(R),
\]
where $S$ acts on the target through $\mu$. Define
$D_e(R/k)=\im\rho_e$. Thus $D_e$ consists exactly of the maps
that admit an actual diagonal-compatible lift.

Let $k'/k$ be an extension of perfect fields, and put
$R'=R\otimes_k k'$. The natural isomorphism
\[
 R'\otimes_{k'}R'=S\otimes_k k'
\]
identifies its multiplication ideal $I'$ with $I\otimes_k k'$.

\begin{theorem}\label{thm:descent}
There are natural module identifications
\[
 H_e(R')=H_e(R)\otimes_k k',\qquad
 C_e(S\otimes_k k',I')=C_e(S,I)\otimes_k k',
\]
which identify the restriction maps. Consequently
\[
 D_e(R'/k')=D_e(R/k)\otimes_k k'.                         \tag{1}
\]
For $h\in H_e(R)$, its extended Cartier map is diagonally liftable
over $R'$ if and only if $h\in D_e(R/k)$.
\end{theorem}

The displayed tensor identifications include the Frobenius twist
specified below. They do not assert that a Cartier map acts $k'$-linearly
on the untwisted additive vector space.

\section{Finiteness and the Frobenius identifications}
Every finite type $k$-algebra $B$ has finite Frobenius. Indeed choose
finitely many algebra generators $x_1,\ldots,x_a$. Every polynomial
in them is a sum of expressions $b_\alpha^q x^\alpha$ with
$0\leq\alpha_j<q$, because coefficients have $q$-th roots in $k$.
Their images in any quotient retain that expression. Thus the finitely
many $F_*^ex^\alpha$ generate $F_*^eB$ as a $B$-module. As $B$ is
Noetherian, this module is finitely presented. For a finite ideal $J$
of $B$, $F_*^eJ$ is a submodule of this finite module, hence is finite
and finitely presented as well. This applies to $R,S,I$ and their
field extensions.

For any such $B$, let $B'=B\otimes_k k'$. There is a $B'$-linear
isomorphism
\[
 \Phi_B:(F_*^eB)\otimes_k k'\longrightarrow F_*^eB',\qquad
 (F_*^ex)\otimes c\longmapsto F_*^e(x\otimes c^q).          \tag{2}
\]
The tensor on the left uses the $k$-module structure inherited from
the Frobenius $B$-module. In particular
$(F_*^e(\lambda^q x))\otimes c=(F_*^ex)\otimes\lambda c$.
This relation is preserved by (2), since
$\lambda^q x\otimes c^q=x\otimes(\lambda c)^q$.
Additivity holds because raising to $q$ is additive in characteristic
$p$. The inverse sends $F_*^e(x\otimes c)$ to
$(F_*^ex)\otimes c^{1/q}$ and extends additively. It is well defined:
the relation $\lambda x\otimes c=x\otimes\lambda c$ becomes the
same tensor relation using $\lambda^{1/q}\in k$.
Perfectness of both fields supplies all these roots uniquely.

For $a\in B$ and $b,c\in k'$, the action of $a\otimes b$ on the
right gives $F_*^e(a^q x\otimes b^q c^q)$; this is also the image of
$a(F_*^ex)\otimes bc$ on the left. Hence (2) and its inverse are
$B'$-linear. The same formula holds for $J$ in place of $B$ and
identifies the inclusions $F_*^eJ\subset F_*^eB$ and their quotients
with the extended inclusions and quotients. No embedding of $B$ in
a ring of $q$-th roots, or assumption that $B$ is reduced, is used.

\section{Hom modules, compatibility kernels, and restriction images}
If $M$ is a finitely presented module over a ring $B$ and $B'$ is flat
over $B$, then
\[
 \Hom_B(M,N)\otimes_B B'
    =\Hom_{B'}(M\otimes_B B',N\otimes_B B').              \tag{3}
\]
To prove this directly choose a finite presentation
$B^a\to B^b\to M\to0$. The Hom module is the kernel of
$N^b\to N^a$. Flat tensoring preserves that kernel and the finite
direct sums, giving (3). This proof requires no flatness of $N$.

Apply (3) to $F_*^eR$ and $F_*^eS$, using (2) for their extended
domains. Tensoring by $B\otimes_k k'$ over $B$ is the same as
tensoring by $k'$ over $k$. It gives the asserted identification of
$H_e$ and of the ambient Hom modules on $S$.
The compatibility module is precisely the kernel of
\[
 \Hom_S(F_*^eS,S)\longrightarrow\Hom_S(F_*^eI,S/I),\qquad
 \psi\longmapsto(\psi|_{F_*^eI}\bmod I).                 \tag{4}
\]
The finite presentation of $F_*^eI$ permits (3) on the target also.
The identifications (2) respect restriction and reduction modulo $I$,
so (4) after flat extension is exactly the compatibility kernel for
$(S\otimes_k k',I')$. Flatness preserves its kernel, proving the
asserted identification of $C_e$.

For a compatible $\psi$, define $\rho_e(\psi)$ by choosing a lift
of an element of $F_*^eR$ to $F_*^eS$, applying $\psi$, and reducing
modulo $I$. Different lifts differ by $F_*^eI$ and therefore give
the same result. The map is $R$-linear for the Frobenius source action.
The functorial identifications (2) commute with $S\to R$, and hence
the extended restriction is exactly $\rho'_e$. Taking the image of
this map commutes with flat tensor product. This proves (1).

Finally form the cokernel $Q=H_e(R)/D_e(R/k)$. Its field extension
is $H_e(R')/D_e(R'/k')$ by flatness. For an element $h$, its class
is zero after extension if and only if it was zero beforehand, since
$k\to k'$ is faithfully flat. This proves the liftability assertion
and completes the theorem.

\begin{corollary}\label{cor:ideal}
Let $J\subset R$ be an ideal. If
$J R'\cdot H_e(R')\subset D_e(R'/k')$, then
$J\cdot H_e(R)\subset D_e(R/k)$.
\end{corollary}
\begin{proof}
For $a\in J$ and $h\in H_e(R)$, the extension of $ah$ is
$(a\otimes1)(h\otimes1)$. By hypothesis it belongs to the extended
image. The final assertion of Theorem~\ref{thm:descent} gives
$ah\in D_e(R/k)$. Sums of these elements prove the inclusion.
\end{proof}

\section{Exact finite-field use and limitations}
When $k$ is finite, take the infinite perfect extension
$k'=\bigcup_{r\geq0}k(t^{1/p^r})$. The preceding proof applies:
$R'$ and $S'$ are still finite type algebras over the perfect field
$k'$, hence Noetherian and Frobenius-finite. No finiteness of $k'/k$
is required. No multiplicative field retraction $k'\to k$ is used.

For a fixed index $d$, differentials and a finite presentation show
\[
 \Omega_{R'/k'}=\Omega_{R/k}\otimes_k k',\qquad
 \Fitt_d(\Omega_{R'/k'})=\Fitt_d(\Omega_{R/k})R'.          \tag{5}
\]
The second equality follows by extending the matrix and its minors.
Thus Corollary~\ref{cor:ideal} applies to the Jacobian ideal whenever
it is defined by this same index on the two algebras. In an application
to reduced equidimensional finite type algebras, one must additionally
check the retained dimension and geometric-reducedness hypotheses
needed by a separately applied infinite-field theorem.

Smolkin, \emph{A new subadditivity formula for test ideals},
\href{https://arxiv.org/abs/1805.08739v3}{arXiv:1805.08739v3},
Theorem 3.12, seeks the inclusion
$\operatorname{Jac}(R/k) H_e(R)\subset D_e(R/k)$.
Its printed finite-field step regards $k$ as an $L$-module via a
$k$-linear retraction $L\to k$, then tensors an $L$-linear diagram
by $k$. A merely linear retraction does not provide that module
structure. Corollary~\ref{cor:ideal}, with (5), supplies a valid
finite-field reduction \emph{provided the separately justified
infinite-perfect-field inclusion and its ambient hypotheses hold}.
It is a descent of the actual compatible Hom kernel and restriction
image, not a descent of test ideals via a guessed equality.

The full printed statement and immediate proof of Smolkin 3.12,
including both the problematic finite-field paragraph and its separate
infinite-field construction, have been read. The relevant printed
page was physically checked earlier. The infinite-field argument's
Hochster--Huneke Noether-normalization input is a separate published
foundation; this paper does not claim a fresh reconstruction of that
input or grant correctness to the entire source proof. The module
theorem above was proved independently and does not require it.

Neither (1) nor (5) proves a test-ideal base-change identity,
singular subadditivity, a trace-image formula, or a uniform all-support
envelope estimate. Those require their own exact statements and proofs.
The original normal excellent local-envelope problem P12.3, including
nonisolated, mixed-characteristic and nonalgebraic residuals, remains
open. This working paper supplies no new original closure or accepted
paper count. It does not alter frozen paper377 or provide feedback to
the active independent campaign269. Full primary bodies and pixels
remain local research data, and no supplied program was executed.
\end{document}
